% =============================================================================
% introduction.tex -- Introduction, written for a reader with no prior exposure to
% membership inference or memorization measurement.
%
% REVISION CONTRACT (advisor feedback, 2026-07):
% * No technical term is used here before it is explained. Statistical machinery
%   (partial correlation, mediation, collinearity, low-FPR evaluation) is NAMED but
%   its definition is deferred to Section 2, which is signposted explicitly.
% * The contamination -> memorization -> leakage chain is presented as an ARGUMENT
%   WITH ASSUMPTIONS, not as an assertion. The assumptions are enumerated in
%   Section 2.4 and referenced here.
% * Scope: we measure EXTRACTION. Sensitive-information leakage is an assumption of
%   the framing, not a result of this paper. Our PII measurement returned a null at
%   this scale and is reported as such.
% =============================================================================

\section{Introduction}
\label{sec:intro}

\subsection{Benchmarks, and the assumption they depend on}
Large language models (LLMs) are compared, selected, and increasingly certified as safe
on the basis of their scores on public test sets called
\emph{benchmarks}~\citep{hendrycks2021mmlu,cobbe2021gsm8k}. A benchmark is a fixed
collection of questions with known answers. A model that answers more of them is taken
to be more capable.

That inference depends on one assumption: the model must not have seen the benchmark
during training. If it has, a high score may reflect recall of the answer rather than
the ability the benchmark was built to measure. The presence of evaluation data inside a
model's training corpus is called \emph{benchmark
contamination}~\citep{golchin2024timetravel}.

Contamination is now difficult to avoid. Benchmarks are small, fixed, and copied widely
across the web once published. Training corpora are enormous and lightly filtered,
assembled at the scale of hundreds of gigabytes to
petabytes~\citep{commoncrawl,gao2020pile}. Benchmark items are therefore drawn into the
next crawl through ordinary redistribution, with no adversary and no misconduct
required. The research community treats this primarily as a measurement-hygiene
problem, on the grounds that a contaminated score overstates
capability~\citep{ravaut2024survey}.

\subsection{Why this is also a privacy and security question}
There is a second consequence that the hygiene framing sets aside. The mechanism that
makes a contaminated benchmark score untrustworthy is \emph{memorization}: the model has
retained information specific to individual training examples rather than only the
general patterns across them. Memorization is also the mechanism behind a known privacy
failure. Models trained on web-scale corpora can reproduce training text word for word,
including personal information such as names, addresses, and contact
details~\citep{carlini2021extracting,carlini2023quantifying}.

This suggests an appealing argument: because contamination is easy to look for and
sensitive-data leakage is not, perhaps a cheap contamination signal can act as an early
warning for leakage risk. If so, contaminating a benchmark would not merely distort a
leaderboard. It would expose a measurable channel into the training data.

We want to be careful here, because this argument is not automatic, and stating it
loosely has consequences for what the resulting measurements mean. Contamination and
sensitive-data leakage are related but genuinely distinct phenomena. Contamination
concerns \emph{evaluation} data appearing in a training corpus. Leakage concerns a model
\emph{emitting} sensitive content. One does not imply the other. Benchmark items are in
fact usually public and non-sensitive by construction, so a contaminated benchmark is
not itself a privacy harm. Any bridge between the two must be built from explicit
premises about a shared underlying mechanism.

Section~\ref{sec:chain} therefore sets out that bridge as four named assumptions
(A1 through A4) rather than as a claim, and states plainly which of them this paper
tests, which it inherits from prior work, and which remain untested. The short version
is that our experiments bear on the first link, whether a contamination score reflects
item-specific retention at all, and that the link from retention to \emph{sensitive}
content remains an assumption of the framing rather than a finding of this paper.

\subsection{The practical problem: which instrument should an auditor trust?}
Suppose you want to know whether a model has memorized a particular document. The
direct test is to try to make the model reproduce it, which requires generating text and
comparing it against the original. That is expensive at corpus scale and requires
possessing the document already.

The cheaper alternative, and the one widely used in practice, is a \emph{detector}: a
score computed from the model's own reported probabilities for the text, with no
retraining and no reference models. Several such detectors exist, and Section
\ref{sec:concepts} defines each one before we use it. The simplest is the model's
\emph{loss} on the text, which measures how unsurprised the model is by it. Others
(Min-K\%, Min-K\%++, and a compression-based ratio) are transformations of the same
per-token probabilities, designed to separate training members from non-members more
sharply than loss alone.

These detectors are developed and ranked by how well they answer a membership question:
was this text in the training data? But an auditor concerned with privacy wants the
answer to a different question: how much of this specific item has the model retained,
and would it come back out? Those two questions need not have the same answer, and the
distinction has been noted before. \citet{hayes2025strong} report no correlation between
a strong membership attack's success and extraction on the models they study, and
suggest the two capture different signals.

Our question is the one that follows from theirs, and it is a question about measurement
validity rather than about attack strength. If the detectors an auditor would actually
reach for are all computed from the same per-token probabilities as loss, do they carry
any information about leakage that loss does not already carry? This matters because it
determines whether detector-based auditing can be trusted to tell you \emph{which} items
are at risk, or only whether contamination happened at all.

\subsection{Research questions}
\begin{itemize}
 \item \textbf{RQ1 (conceptual validity).} Under what assumptions can a contamination
 score serve as an indicator of memorization, and further as an indicator of
 sensitive-information leakage? Which of those assumptions are empirically testable in a
 setting where training-set membership is known with certainty?
 \item \textbf{RQ2 (incremental predictive value).} Do the widely used detectors provide
 information about which individual items a model has memorized, over and above what the
 model's per-item loss already provides?
 \item \textbf{RQ3 (choice of instrument).} These detectors were designed and tuned to
 answer one question well, namely whether a text was in the training data. Do the design
 adjustments \emph{intended} to sharpen that discrimination also help predict whether the
 model will reproduce the text, or do the two objectives come apart? We note in advance
 that at our scale no detector separates members from non-members above chance
 (Section~\ref{sec:res-membership}), so RQ3 is answered here only in the weaker,
 design-level sense.
\end{itemize}

\subsection{Approach in brief}
We study Pythia models~\citep{biderman2023pythia}, which were trained on The
Pile~\citep{gao2020pile}, a corpus that is public. This is the central design choice of
the paper and Section~\ref{sec:eval} defends it in detail. Because the corpus is
published, we can look up whether a document was genuinely in the training data instead
of inferring it. Membership is therefore ground truth, which removes the largest source
of error in this literature, namely that a detector can appear to succeed because the
member and non-member texts differ in topic or era rather than in
membership~\citep{duan2024mia}.

For each document we compute the detector scores, and separately measure a concrete
leakage outcome: given the first part of the document, does the model reproduce the
rest~\citep{carlini2023quantifying}? We then ask whether the detector scores predict
that outcome once the model's loss is held fixed. Section~\ref{sec:concepts} explains
what "holding loss fixed" means and why it is the right comparison, and
Section~\ref{sec:eval} explains why we pre-committed to the analysis in writing before
running it.

\subsection{Contributions, and explicit non-contributions}
We are deliberate about what this paper is not. It introduces no new detector, no new
attack, and no new metric. Every detection method we run is from prior
work~\citep{yeom2018privacy,shi2024detecting,zhang2025minkpp,carlini2021extracting,brown2020gpt3,oren2024proving}.
We do not train or fine-tune models, and we do not attack deployed systems for real
third-party personal data. Within that scope:

\begin{itemize}
 \item \textbf{An explicit conceptual bridge, stated as assumptions.} We replace the
 usual one-line assertion that contamination indicates leakage risk with four named
 assumptions, and we identify which are testable, which we test, and which the field
 currently takes on faith (Section~\ref{sec:chain}). This makes the framing falsifiable
 and bounds what any contamination-based privacy claim can support.
 \item \textbf{A ground-truth comparison of existing detectors as leakage predictors.}
 On Pythia and The Pile, where membership is known, we evaluate the detectors both as
 membership classifiers and as predictors of a per-item extraction outcome, using the
 low-false-positive reporting convention that privacy work
 expects~\citep{carlini2022lira} and controlling for the frequency, duplication, and
 distribution-shift confounds that prior work
 identifies~\citep{biderman2023pythia,duan2024mia}.
 \item \textbf{A controlled measurement of incremental value, which is the paper's core
 result.} Using an analysis fixed in advance, we ask whether each detector predicts
 leakage after the model's per-item loss is accounted for. In our setting the detectors
 add no positive predictive value beyond loss. They are close to being mathematical
 restatements of loss, which we quantify directly, and we therefore read the result as a
 measurement-validity finding: these instruments carry little information about
 \emph{which} items leak that loss does not already carry.
\end{itemize}

\paragraph{What we do not claim.} We do not claim to demonstrate leakage of sensitive or
personal information. We attempted such a measurement on the Enron Emails subset that
sits inside The Pile, and it returned no detected leakage at the scale we could run
(Section~\ref{sec:results}). We report that null honestly and treat the step from
memorization to \emph{sensitive} content as an untested assumption
(A3 and A4 in Section~\ref{sec:chain}) rather than as a result. We also do not claim a
stronger attack than prior work, and we do not claim that the detectors are
\emph{worse} than loss, only that we find no evidence they are better.

\paragraph{Preliminary scope.} The results reported here come from the smallest Pythia
model, at $160$ million parameters, run on a CPU. Memorization increases with model
scale~\citep{carlini2023quantifying}, so this is the regime least likely to show
leakage, and we flag every affected conclusion as preliminary. The pipeline was built so
that repeating everything on larger models is a change of one configuration value.

\subsection{Roadmap}
Section~\ref{sec:background} explains benchmarks, memorization, and the technical terms
and statistical tools this paper relies on, and then states the assumption chain that
connects contamination to leakage. Section~\ref{sec:threat} defines the threat model.
Section~\ref{sec:relatedwork} positions the work against prior contamination,
memorization, and membership-inference research. Section~\ref{sec:eval} gives the
experimental design and the reasoning behind each methodological choice.
Section~\ref{sec:results} reports results, Section~\ref{sec:discussion} interprets them,
and Section~\ref{sec:limitations} states the limitations.
